\documentclass[12pt,a4paper]{article}
\usepackage{geometry}
\geometry{left=2.5cm,right=2.5cm,top=2.0cm,bottom=2.5cm}
\usepackage[english]{babel}
\usepackage{amsmath,amsthm}
\usepackage{amsfonts}
\usepackage[longend,ruled,linesnumbered]{algorithm2e}
\usepackage{fancyhdr}
\usepackage{ctex}
\usepackage{array}
\usepackage{listings}
\usepackage{color}
\usepackage{graphicx}

\begin{document}

\title{
  {\heiti《算法分析与设计》第 {$1$} 次作业
    \footnote{要求：1、分析题请用书面化语言给出详细分析过程。}
    }
}
\date{}

\author{
姓名：\underline{你的名字}~~~~~~
学号：\underline{你的学号}~~~~~~
成绩：\underline{~~~~~~~~~~~~~~~~~~}
}

\maketitle

\noindent
\section*{\bf \color{red}{算法分析题}}

\noindent
{\bf 题目1：}求下列函数或级数和的$\Theta$渐进表达式。
\begin{enumerate}
  \item[(1)]  $F(n)=6n^{6}+6^n$
  \item[(2)]  $R(n)=3n^3 + n^2 \log n^4$
  \item[(3)]  $H(n)=\frac{n (\log n + n \sqrt{n})}{n^2 + \log n}$
  \item[(4)]  $G(n)=\sum\limits_{k=1}^n\frac{k(\sqrt{k} \log k+1)}{\sqrt{k}+1}$
  \item[(5)]  $T(n)=\sum\limits_{k=1}^n\frac{3^k(k^2+5k+3)}{2k(k+5)}$
\end{enumerate}

\vspace{5pt}
\noindent
{\bf 答：}


\vspace{10pt}
\noindent
{\bf 题目2：}用$\Theta$记号表示对下面一段程序中语句$x \gets x+6$被执行的次数的估计。
\begin{algorithm}
\caption{一个算法实例}

\For{int i=1 \emph{\KwTo} n}{
    \For{int j=$\frac{i}{3}$ + 1 \emph{\KwTo} i}{
    $x \gets x+6$\;
    }
	}
\end{algorithm}

\vspace{5pt}
\noindent
{\bf 答：}


\section*{\bf \color{red}{证明题}}

\vspace{10pt}
\noindent
{\bf 题目3：}假设$f(n)$，$g(n)$是两个定义域为自然数的正函数。证明 $f(n)+g(n)=\varTheta(\max(f(n),g(n)))$

\vspace{5pt}
\noindent
{\bf 答：}

\end{document}